\documentclass[12pt]{article}
\usepackage[utf8]{inputenc}
\usepackage[T1]{fontenc}
\usepackage{hyperref}
\title{Yukpa Language Resources for K-12 Curriculum Development}
\author{Compiled by Claude Sonnet 5.5 (claude-sonnet-5-5) for Andr\'as Kornai}
\date{}
\begin{document}
\maketitle

\section{Introduction}
Yukpa (also Yuko, Yupa; ISO 639-3 yup; Glottocode yukp1241) is a Cariban language of the Serran\'ia de Perij\'a in Colombia and Venezuela (Zulia). Figures vary: roughly 3,000 speakers in each country in Ethnologue-based entries, with most adults monolingual and younger people increasingly bilingual in Spanish. The searches found sociolinguistic and curriculum-related reports but no dictionary or primer.

\section{Primary Educational Resources}
\begin{itemize}
\item Colombian Ministry of Education page about a Yukpa curriculum project in the Serran\'ia del Perij\'a, led by a teacher-training school (Spanish): \url{https://www.mineducacion.gov.co/1621/article-208171.html}
\end{itemize}

\section{Linguistic and Reference Materials}
\begin{itemize}
\item Sociolinguistic study of Yukpa (2007--2009, Colombia and Venezuela) covering dialects, population and education, Universidad Nacional de Colombia: \url{https://www.humanas.unal.edu.co/colantropos/files/1014/7948/1900/yukpa.pdf}
\item Endangered Languages Project entry (dialects, vitality): \url{https://www.endangeredlanguages.com/lang/632}
\item Colombian Ministry of Culture profile of the Yukpa people: \url{https://mng.mincultura.gov.co/prensa/noticias/Documents/Poblaciones/PUEBLO%20YUKPA.pdf}
\item A UNAD repository document that appeared in results (contents not verified): \url{https://repository.unad.edu.co/jspui/bitstream/10596/2544/3/37945838.pdf}
\end{itemize}

\section{Religious and Community Resources}
None found.

\section{Digital and Audio Resources}
None found.

\section{Curriculum Development Context}
The Yukpa are reported to teach the language as a subject in some schools, and a Colombian teacher-training curriculum effort exists, but no teaching materials were located. Dialect variation (Rio Casacar\'a, Becerril, and the Venezuelan variety) must be taken into account. Next steps would be to look for Colombian and Venezuelan dictionary projects and university work, which were not found here.

\end{document}